% GENERATED FILE. Edit canonical lessons, then run npm run generate:books.
% canonical-source-hash: fnv1a64:3d260a0d3bf2013d
% canonical-lessons: GE-C399-lenkrad, GE-C399-lenken, GE-C399-bremsen, GE-C399-bremse, GE-C399-schalten

\chapter{Am Steuer}
\label{ch:ge-am-steuer}
\hlchaptermodality{399}

\begin{chapteropening}
\textbf{By the end of this chapter:} \emph{I can describe driving: the steering wheel, steering, braking, the brake and changing gear.}

Complete the last lesson of chapter 399: I can describe driving: the steering wheel, steering, braking, the brake and changing gear.
\end{chapteropening}

\section[das Lenkrad]{das Lenkrad — the steering wheel}

\label{lesson:GE-C399-lenkrad}



\begin{quote}
Before the new one: say the German for the roundabout, then the German for to turn off.
\end{quote}

\subsection*{You'll want to know: das Lenkrad}
\textbf{das Lenkrad} — \textquotedblleft{}the steering wheel\textquotedblright{}.

\begin{quote}
Mein Vater hält das Lenkrad mit beiden Händen.
\end{quote}

\begin{quote}
\emph{My father holds the steering wheel with both hands.}
\end{quote}

Plural: \emph{die Lenkräder}. Also \emph{das Steuer}: \emph{am Steuer sitzen} — to be at the wheel.

\begin{cousinweb}
\emph{Lenkrad} is a compound of \emph{lenken} 'to steer' and \emph{Rad} 'wheel' (EtymDB lists it under \emph{Rad}). Both parts are inherited Germanic words.
\end{cousinweb}

\begin{grammarlens}[title={What you've built}]
One more word for telling, explaining and describing at B1.
\end{grammarlens}

\subsection*{Your turn}
\begin{itemize}
\raggedright
  \item \emph{Say it:} \emph{das Lenkrad}
  \item \emph{Say it:} \emph{das Lenkrad}, once more
  \item \emph{Say it:} a sentence of your own with \emph{das Lenkrad}
  \item \emph{Recall:} say the German for the roundabout, then the German for to turn off, then say \emph{das Lenkrad} again
\end{itemize}

\begin{tcolorbox}[breakable,colback=teal!4,colframe=teal!35!black,title={Before you move on}]
What does das Lenkrad mean? (\textquotedblleft{}The steering wheel\textquotedblright{}.) Say it once more.
\end{tcolorbox}
\section[lenken]{lenken — to steer; to direct}

\label{lesson:GE-C399-lenken}



\begin{quote}
Before the new one: say the German for to turn off, then the German for the steering wheel.
\end{quote}

\subsection*{You'll want to know: lenken}
\textbf{lenken} — \textquotedblleft{}to steer; to direct\textquotedblright{}.

\begin{quote}
Auf dem Eis kann er das Auto kaum lenken.
\end{quote}

\begin{quote}
\emph{On the ice he can hardly steer the car.}
\end{quote}

Also figurative: \emph{das Gespräch auf ein Thema lenken} — to steer the talk towards a topic.

\begin{cousinweb}
\emph{lenken} is an inherited Germanic verb that first meant 'to bend, to turn' (EtymDB gives 'to turn'). EtymDB links it with English \emph{lank} 'limp, thin'.
\end{cousinweb}

\begin{grammarlens}[title={What you've built}]
One more word for telling, explaining and describing at B1.
\end{grammarlens}

\subsection*{Your turn}
\begin{itemize}
\raggedright
  \item \emph{Say it:} \emph{lenken}
  \item \emph{Say it:} \emph{lenken}, once more
  \item \emph{Say it:} a sentence of your own with \emph{lenken}
  \item \emph{Recall:} say the German for to turn off, then the German for the steering wheel, then say \emph{lenken} again
\end{itemize}

\begin{tcolorbox}[breakable,colback=teal!4,colframe=teal!35!black,title={Before you move on}]
What does lenken mean? (\textquotedblleft{}To steer; to direct\textquotedblright{}.) Say it once more.
\end{tcolorbox}
\section[bremsen]{bremsen — to brake; to slow down}

\label{lesson:GE-C399-bremsen}



\begin{quote}
Before the new one: say the German for the steering wheel, then the German for to steer.
\end{quote}

\subsection*{You'll want to know: bremsen}
\textbf{bremsen} — \textquotedblleft{}to brake; to slow down\textquotedblright{}.

\begin{quote}
Der Fahrer bremst, weil ein Hund über die Straße läuft.
\end{quote}

\begin{quote}
\emph{The driver brakes because a dog runs across the road.}
\end{quote}

\emph{Er bremst}, \emph{er hat gebremst}. Figurative: \emph{jemanden bremsen} — to hold someone back.

\begin{cousinweb}
\emph{bremsen} is built on the noun \emph{Bremse} 'brake', which is usually traced to Middle Low German \emph{premse} 'a clamp'. EtymDB files the verb under its look-alike \emph{Bremse} 'horsefly', so the line is hedged.
\end{cousinweb}

\begin{grammarlens}[title={What you've built}]
One more word for telling, explaining and describing at B1.
\end{grammarlens}

\subsection*{Your turn}
\begin{itemize}
\raggedright
  \item \emph{Say it:} \emph{bremsen}
  \item \emph{Say it:} \emph{bremsen}, once more
  \item \emph{Say it:} a sentence of your own with \emph{bremsen}
  \item \emph{Recall:} say the German for the steering wheel, then the German for to steer, then say \emph{bremsen} again
\end{itemize}

\begin{tcolorbox}[breakable,colback=teal!4,colframe=teal!35!black,title={Before you move on}]
What does bremsen mean? (\textquotedblleft{}To brake; to slow down\textquotedblright{}.) Say it once more.
\end{tcolorbox}
\section[die Bremse]{die Bremse — the brake}

\label{lesson:GE-C399-bremse}



\begin{quote}
Before the new one: say the German for to steer, then the German for to brake.
\end{quote}

\subsection*{You'll want to know: die Bremse}
\textbf{die Bremse} — \textquotedblleft{}the brake\textquotedblright{}.

\begin{quote}
Die Bremsen an meinem Fahrrad sind kaputt.
\end{quote}

\begin{quote}
\emph{The brakes on my bike are broken.}
\end{quote}

Plural: \emph{die Bremsen}. \emph{Die Bremse} is also a horsefly, a different word that only looks the same.

\begin{cousinweb}
\emph{Die Bremse} 'brake' is usually traced to Middle Low German \emph{premse} 'a clamp', first one put on a restless horse's nose. EtymDB only lists the look-alike \emph{Bremse} 'horsefly', inherited from Middle Low and Middle High German, so the line for the brake is hedged.
\end{cousinweb}

\begin{grammarlens}[title={What you've built}]
One more word for telling, explaining and describing at B1.
\end{grammarlens}

\subsection*{Your turn}
\begin{itemize}
\raggedright
  \item \emph{Say it:} \emph{die Bremse}
  \item \emph{Say it:} \emph{die Bremse}, once more
  \item \emph{Say it:} a sentence of your own with \emph{die Bremse}
  \item \emph{Recall:} say the German for to steer, then the German for to brake, then say \emph{die Bremse} again
\end{itemize}

\begin{tcolorbox}[breakable,colback=teal!4,colframe=teal!35!black,title={Before you move on}]
What does die Bremse mean? (\textquotedblleft{}The brake\textquotedblright{}.) Say it once more.
\end{tcolorbox}
\section[schalten]{schalten — to change gear; to switch}

\label{lesson:GE-C399-schalten}



\begin{quote}
Before the new one: say the German for to brake, then the German for the brake.
\end{quote}

\subsection*{You'll want to know: schalten}
\textbf{schalten} — \textquotedblleft{}to change gear; to switch\textquotedblright{}.

\begin{quote}
Am Berg schaltet sie in den zweiten Gang.
\end{quote}

\begin{quote}
\emph{On the hill she changes into second gear.}
\end{quote}

\emph{Schalten} also means to switch: \emph{das Licht einschalten} — to switch the light on.

\begin{cousinweb}
\emph{schalten} is inherited through Middle High German \emph{schalten} from Old High German \emph{scaltan} (EtymDB), usually glossed 'to push, to shove'; it is said to have first meant pushing a boat along with a pole, later a lever pushed into place.
\end{cousinweb}

\begin{grammarlens}[title={What you've built}]
One more word for telling, explaining and describing at B1.
\end{grammarlens}

\subsection*{Your turn}
\begin{itemize}
\raggedright
  \item \emph{Say it:} \emph{schalten}
  \item \emph{Say it:} \emph{schalten}, once more
  \item \emph{Say it:} a sentence of your own with \emph{schalten}
  \item \emph{Recall:} say the German for to brake, then the German for the brake, then say \emph{schalten} again
\end{itemize}

\begin{tcolorbox}[breakable,colback=teal!4,colframe=teal!35!black,title={Before you move on}]
What does schalten mean? (\textquotedblleft{}To change gear; to switch\textquotedblright{}.) Say it once more.
\end{tcolorbox}
